\chapter{Conclusiones}
\label{chap:conclusiones}

Tras el análisis de los fundamentos teóricos y la estructuración de los ensayos experimentales propuestos para
este trabajo práctico, se concluye que los objetivos planteados al inicio del informe se cumplen satisfactoriamente.
Los métodos implementados permiten una caracterización apropiada de los amplificadores de baja frecuencia.

El método de reducción a la mitad por sustitución demostró ser una técnica sumamente práctica y directa para determinar
las impedancias de entrada ($Z_i$) y salida ($Z_o$) en régimen sinusoidal de audiofrecuencia. Esta metodología resulta
muy efectiva siempre que la parte reactiva de las impedancias bajo estudio sea despreciable frente a su componente resistiva.

Por otro lado, la utilización de escalas en decibeles ($dB$, $dBu$ y $dBm$)
evidenció su utilidad: se agiliza la determinación de ganancia total en sistemas multietapa, permitiendo
sustituir operaciones de multiplicación por simples sumas y restas de niveles logarítmicos. Además, la escala
logarítmica acomoda de manera natural rangos dinámicos extremadamente amplios de
señales.

El estudio del efecto de la realimentación negativa permitió observar de manera
clara el comportamiento estudiado y por lo tanto, esperado, del sistema
realimentado:
\begin{itemize}

    \item \textit{Estabilización del sistema:} Se pudo observar que la ganancia
        a lazo cerrado disminuye de manera sustancial respecto de la ganancia a
        lazo abierto, lo que se corresponde al comportamiento esperado. Se sabe
        que a pesar de la disminución de ganancia, el circuito gana una enorme
        estabilidad frente a variaciones en los parámetros del transistor,
        fluctuaciones de temperatura y cambios en la tensión de alimentación.
    \item \textit{Modificación de impedancias:} La realimentación de muestra de tensión
        y mezcla serie disminuye drásticamente la resistencia de salida del sistema, haciendo
        que el amplificador se aproxime al comportamiento de una fuente de
        tensión ideal. Al mismo tiempo, incrementa su resistencia de entrada,
        optimizando la adaptación de tensión frente a fuentes de señal
        precedentes.
    \item \textit{Extensión de ancho de banda:} Se constató la relación
        de compromiso ganancia-ancho de banda. Al realimentar negativamente el
        amplificador, se expande significativamente el ancho de banda ($AB$),
        desplazando la frecuencia de corte superior hacia valores más elevados y
        la de corte inferior hacia valores más bajos.
    \end{itemize}

Finalmente, la comparación entre el método de barrido de frecuencia senoidal y el método temporal mediante onda cuadrada puso de manifiesto
que el análisis de respuesta al escalón es una herramienta rápida y potente para estimar el ancho de banda del amplificador. A pesar
de basarse en modelos simplificados, el desvío experimental entre ambos métodos suele ser acotado, validando su
uso en diagnósticos rápidos de laboratorio.

